%% ===========================================================================
%% MASTER CV - Marcus D. Yang
%% ---------------------------------------------------------------------------
%% This file is the FROZEN FORMATTING STANDARD and the FACT SOURCE for every
%% tailored CV. Two rules:
%%   1. The preamble below (geometry, font, \section rule, \erow macro, list
%%      spacing) is frozen. A tailored CV COPIES it verbatim - never retune it.
%%   2. Everything between \begin{document} and \end{document} is CONTENT and
%%      is fully re-orderable / re-groupable per the job description. See
%%      .claude/skills/job-application-assistant/05-cv-templates.md.
%% Compile: lualatex -interaction=nonstopmode main_example.tex   (MUST be 1 page)
%% Visual reference: cv/Marcus Yang Master Resume.pdf
%% ===========================================================================
\documentclass[11pt,letterpaper]{article}

% --- Fonts: Times New Roman clone, bundled with every modern TeX dist -------
\usepackage{fontspec}
\setmainfont{TeX Gyre Termes}
\usepackage{microtype}
\linespread{0.93}   % match the master's tight (Word "single") leading

% --- Page geometry (master: 0.8in sides, ~0.5in top) ----------------------
\usepackage[top=0.5in,bottom=0.5in,left=0.8in,right=0.8in]{geometry}

% --- Section headings: bold, full-width rule underneath --------------------
\usepackage{titlesec}
\titleformat{\section}{\fontsize{12.5}{14}\selectfont\bfseries}{}{0pt}{}[{\vspace{1pt}\hrule height 0.6pt}]
\titlespacing*{\section}{0pt}{6pt}{2pt}

% --- Tight bullet lists ---------------------------------------------------
\usepackage{enumitem}
\setlist[itemize]{leftmargin=1.35em,itemsep=1pt,topsep=1pt,parsep=0pt,
                  partopsep=0pt,label=\textbullet}

\usepackage[hidelinks]{hyperref}
\pagestyle{empty}
\setlength{\parindent}{0pt}

% --- \erow{left}{right}: entry/degree row, left flush-left, right flush-right
%     Caller supplies ALL markup in {left} (\textbf{...}, \textit{...}, etc.).
\newlength{\erowdate}\setlength{\erowdate}{2.0in}
\newcommand{\erow}[2]{%
  \noindent\makebox[\dimexpr\linewidth-\erowdate\relax][l]{#1}%
  \makebox[\erowdate][r]{#2}\par}

% --- \entrygap: vertical space between entries within a section -----------
\newcommand{\entrygap}{\vspace{2pt}}

\begin{document}

% ===========================================================================
%  HEADER
% ===========================================================================
\begin{center}
  {\fontsize{17}{20}\selectfont\bfseries Marcus D. Yang}\\[2pt]
  {\footnotesize (650) 954-0288 \textbar{} \href{mailto:yangdmarcus@gmail.com}{yangdmarcus@gmail.com} \textbar{} \href{https://www.linkedin.com/in/marcusdyang}{www.linkedin.com/in/marcusdyang} \textbar{} \href{https://projects-website-psi.vercel.app/}{https://projects-website-psi.vercel.app/}}
\end{center}

% ===========================================================================
%  EDUCATION
% ===========================================================================
\section{Education}
\erow{\textbf{University of Illinois Urbana-Champaign} \textbar{} Champaign, IL}{August 2024 - May 2028}
\erow{\textit{B.S. in Aerospace Engineering with a Computer Science minor, James Scholar Honors}}{GPA: 3.57/4}
\textbf{Relevant Coursework:} Electrical and Electronic Circuits, Aerospace Control Systems, Discrete Structures, Intro to Computer Science II, Linear Algebra, Dynamics, Aerospace Structures, Thermodynamics, Engineering Materials

% ===========================================================================
%  ORGANIZATIONAL AND TECHNICAL EXPERIENCE
% ===========================================================================
\section{Organizational and Technical Experience}

\erow{\textbf{TVC Rocket} (Individual Project)}{Summer 2026}
\begin{itemize}
  \item Wrote C++ flight firmware on the Teensy 4.1, reading the MPU-6050 IMU over I2C and fusing its data with a complementary filter to drive a closed loop PID gimbal controller.
  \item Designed and fabricated a 2-axis TVC gimbal, avionics bay, and recovery system by integrating a Teensy 4.1 flight computer, IMU, and servo actuators into a commercial rocket body.
  \item Bench tested the PID firmware hardware-in-the-loop against a live IMU feed to confirm gimbal servo response before static fire.
  \item Logged 83 Hz telemetry with zero dropouts over two static fires (15,728 and 12,611 samples) for data analysis.
  \item Root-caused two static fire test failures by analyzing telemetry logs and frame-by-frame video, isolating a misreferenced attitude axis and an inverted control loop sign.
\end{itemize}
\entrygap

\erow{\textbf{Liquid Rocketry at Illinois} (Project Team)}{Spring 2025 - Present}
\begin{itemize}
  \item Designed a KiCad schematic for an RP2040 LoRa telemetry transmitter with USB-C and QSPI flash.
  \item Modeled P\&ID system architecture in Siemens NX and integrated components into the pressurant tank assembly, sourcing and constraining CAD hardware to spatial requirements.
  \item Ran a pressurant tank trade study and produced mechanical integration models, verifying fit against launch vehicle structural constraints and P\&ID.
  \item Created CAD geometries for canard-grid fins and ran Ansys Fluent CFD cases across 0--20\textdegree{} deflection.
\end{itemize}
\entrygap

\erow{\textbf{CS 128/225 Coursework Projects}}{Spring 2025}
\begin{itemize}
  \item Built a C++ database engine with manual memory management, Rule-of-Three semantics, and indexed rows.
  \item Implemented a doubly linked list with a custom iterator and a from-scratch merge sort in C++.
  \item Wrote C++ algorithms for DNA-sequence matching, greedy grid pathfinding, and seam-carving image resizing.
\end{itemize}
\entrygap

\erow{\textbf{Model Rocket Project}}{Fall 2024}
\begin{itemize}
  \item Designed custom fin geometry in Siemens NX and fabricated a physical prototype to hit a target apogee.
  \item Validated flight test data against OpenRocket simulations; flight reached 430 ft apogee with 25 ft downrange drift.
\end{itemize}

% ===========================================================================
%  TECHNICAL SKILLS
% ===========================================================================
\section{Technical Skills}
\textbf{Computer Languages:} Python, C++, Java, JavaScript\\
\textbf{Software:} Siemens NX, Fusion 360, Eagle PCB, KiCad, OpenRocket, XFLR5, Ansys Fluent, Excel\\
\textbf{Shop Tools:} Bambu A1 Mini, Spot Welder, Belt Sander, Laser Cutter, CNC Plasma Cutter, Table Saw, Drill Press\\
\textbf{Materials:} Epoxy Resin, Fiberglass, Sheet Metal, Aluminum Pipes, Plywood, PETG, PLA

% ===========================================================================
%  LEADERSHIP AND ACTIVITIES
% ===========================================================================
\section{Leadership and Activities}

\erow{\textbf{Asian American Association}}{Fall 2024 - Present}
Treasurer \textbar{} Fundraising Officer
\begin{itemize}
  \item Managed membership operations for 300+ members, coordinating records and organizational logistics.
  \item Secured partnerships with 15 on-campus sponsors to support events and fundraising initiatives.
\end{itemize}
\entrygap

\erow{\textbf{Flight Club Aerospace}}{Summer 2021 - Spring 2024}
Landing Gear Research Lead \textbar{} CAD \textbar{} Fabrication
\begin{itemize}
  \item Manufactured wing and control surfaces: cut foam ribs, trimmed aluminum spars with a circular saw, and applied epoxy \& fiberglass layups to shape aerodynamic skins; aligned two 7-ft wings for flight testing.
  \item Drafted gussets and ribs in Fusion 360 and confirmed fit-checks between structural and aerodynamic components.
  \item Oversaw logistics for wing transport/disassembly/reassembly and assembled and presented landing-gear recommendations based on structural load and cost analyses.
\end{itemize}

\end{document}
